%% Section 7 — Discussion
\section{Discussion}
\label{sec:discussion}

\subsection{Multimodality as theoretical necessity}

The paper's methodological argument — that studying datacentering requires combining text, image, financial, and temporal data streams — is sometimes read as a claim about comprehensiveness: more data modalities produce a fuller picture. The argument of this paper is different. Multimodality is not an empirical strategy for reducing gaps but a theoretical necessity following from the CAS structure of the phenomenon.

Each data modality addresses a dimension of the assemblage that the others structurally cannot reach. The corporate text corpus captures the imaginary production through which operators constitute their infrastructure as placeless and their practices as responsible — a communicative act that shapes what becomes legible as a planning problem. Satellite imagery captures the material footprint independently of operator disclosure, providing the state-derived account against which imaginary claims can be read. Financial and ownership data reach the legal and investment architecture through which accountability is distributed or dissolved — a dimension entirely invisible to both text and imagery. And contestation signal data captures the temporal dynamics through which community mobilisation unfolds, escalates, and either produces regime response or dissipates without one.

The opacity of data center infrastructure is not a data gap to be solved with more comprehensive collection. It is a sociotechnical fact: the deliberate, sustained construction of illegibility is itself a constitutive feature of the assemblage, produced through architectural choices, legal structures, communications strategies, and the ``cloud'' metaphor's ideological work. A firm that will not publish its water consumption can still be found in a REIT filing. A campus designed to be visually indistinguishable from a warehouse is still visible in a Sentinel-2 image. A planning moratorium leaves traces in municipal council records even when operators prefer it unexamined. Multimodality is the methodological response to deliberate opacity: triangulation around the assemblage, not penetration of it.

\subsection{The open exchange as theoretically motivated design}

The most consequential implication of the CAS framing for research design is the one most likely to be read as peripheral to the scholarly contribution: the open research infrastructure. The argument of this paper is that the open exchange is not a nice-to-have or an act of scholarly generosity, but a theoretically motivated design choice that follows from the properties of the phenomenon.

Data center assemblages span hundreds of jurisdictions, thousands of facilities, multiple layers of ownership obscured by deliberate legal complexity, and contestation dynamics unfolding across years and decades. The knowledge required to map this — local expertise about specific planning disputes, access to national company registries, linguistic capacity across the jurisdictions involved, methodological knowledge across NLP, computer vision, network analysis, and event detection — is distributed across researchers, civil society organisations, journalists, activists, and community members in ways that no single research team can replicate. A CAS-like phenomenon of this scale requires a distributed observational infrastructure to match.

The research commons — the community-maintained facility index, the shared analytical pipelines, and the open data — is the governance architecture of a distributed observatory. The contributing guide that accompanies this paper is not supplementary documentation; it is part of the methodological argument, specifying the protocols through which distributed knowledge can be aggregated, quality-controlled, and made analytically tractable. When the bibliography of this paper includes a contribution from a local planner in Ashburn who submitted a facility record, or a journalist in Singapore who documented an unreported data center campus, the knowledge produced will be qualitatively different from what any single research team could produce — not just broader, but differently constituted, with a different relationship to the phenomena it describes.

This is the sense in which the open exchange is theoretically necessary: not because collaboration is virtuous, but because the CAS structure of the object demands a research architecture capable of matching its scale, distribution, and deliberate opacity.

\subsection{Positioning within innovation studies}

The paper's argument, if accepted, has implications for the broader innovation studies research programme that extend beyond the specific empirical domain of data centers.

The first implication concerns the object. Data center infrastructure is not a niche concern for digital geographers or infrastructure studies specialists. It is a transition-critical technology: the material substrate through which energy transitions, platform governance, digital sovereignty, and the greening of the economy are being simultaneously negotiated. An electricity consumer of 400~TWh annually that is growing at double-digit rates cannot remain outside the sustainability transitions research programme without that programme sustaining a significant blind spot in its understanding of the socio-technical regime within which energy transitions are occurring.

The second implication concerns the concept. ``Datacentering'' as a gerund is a methodological invitation as much as a theoretical claim. It asks researchers to shift their attention from the stable properties of a technology to the ongoing, contested, processual work of assembling, reproducing, and resisting it. This shift — from noun to gerund, from object to process — is the move that makes the MLP's distinction between landscape, regime, and niche applicable: the ``cloud'' regime is being assembled and contested in real time, and the contestation dynamics documented in this paper are the niche-level disruptions through which transition pressure is beginning to accumulate.

The third implication concerns the CAS framing itself. The risk of invoking complexity theory is that it becomes a theoretical free pass — a way of gesturing at systemic properties without specifying what they are or how they would be detected. The Phillips and Ritala (2019) three-dimensional structure disciplines this application: by committing to the conceptual, structural, and temporal dimensions as distinct but interrelated analytical targets, the paper specifies what an adequate account of the assemblage must reach, and it provides the methodological apparatus to reach each dimension. This is the sense in which the CAS lens functions as sensemaking rather than metaphor: it organises the research design without overclaiming about the formal properties of the system under study.
